\section{Referencial Teórico}
\label{referencial}

Esta seção apresenta os conceitos que fundamentam o estudo e a evolução das ideias até o estado da arte. Ela percorre a passagem da Engenharia de Software tradicional aos LLMs, os paradigmas ES 2.0 e ES 3.0, os agentes inteligentes, os sistemas multiagentes, a orquestração, as capacidades de agentes baseados em LLMs, a programação baseada em agentes e a integridade da avaliação. Ao final, a Subseção~\ref{ssec:correlatos} discute os trabalhos correlatos e posiciona este estudo em relação a eles.

\subsection{Engenharia de Software: da Crise do Software aos LLMs}
\label{ssec:origens}

A Engenharia de Software nasceu de uma crise. Nas décadas de 1960 e 1970, a complexidade crescente dos sistemas expôs os limites da programação informal: orçamentos estourados, prazos descumpridos e projetos abandonados. A conferência da NATO de 1968 \citep{Naur1969NATO} reconheceu o problema e propôs tratar o desenvolvimento de software com o rigor das demais engenharias. As respostas vieram em ciclos: o modelo em cascata sistematizou as fases do processo; os métodos ágeis, consolidados no Manifesto Ágil de 2001 \citep{Beck2001AgileManifesto}, introduziram ciclos iterativos e incrementais; e o DevOps integrou desenvolvimento e operação em \textit{pipelines} de integração e entrega contínuas. Em todos esses ciclos, o desenvolvedor humano permaneceu no centro do processo: ele interpretava requisitos, decidia e escrevia o código.

A Inteligência Artificial entrou no processo primeiro como apoio pontual, em detecção de defeitos, sugestão de refatorações e priorização de testes. A mudança de escala veio com os LLMs, modelos treinados sobre grandes volumes de texto e código, capazes de gerar linguagem natural e código-fonte a partir de instruções. Com eles, a IA deixou de apoiar tarefas isoladas e passou a participar do desenvolvimento, o que deu origem aos paradigmas descritos a seguir.

\subsection{Engenharia de Software 2.0 e 3.0}
\label{ssec:es2es3}

Na Engenharia de Software 2.0, engenheiros humanos usam ferramentas de IA generativa como copilotos \citep{ES302024}. Ferramentas como o GitHub Copilot \citep{GitHubCopilot2024}, o Cursor \citep{Cursor2022} e o Claude Code \citep{ClaudeCode2025} aceleram codificação, depuração e geração de testes. O fluxo decisório, porém, continua com o desenvolvedor, que controla o processo e responde pelos artefatos. \cite{ES302024} apontam três limites desse arranjo. O desenvolvedor sofre sobrecarga cognitiva, porque permanece no centro do laço de criação do código: decompõe o problema, instrui o assistente, avalia as sugestões, depura as falhas e itera. O treino dos modelos é ineficiente e pouco eficaz. E o código gerado tem qualidade subótima e um viés aditivo, que o torna mais difícil de manter; os autores registram ainda que edições em muitos arquivos continuam pouco confiáveis.

Em resposta a esses limites, \cite{ES302024} propuseram a Engenharia de Software 3.0, um paradigma centrado em intenção e orientado a conversa, no qual a IA atua como parceira de equipe e não apenas como assistente. Os autores propõem uma pilha tecnológica conceitual composta por:
% Descrições conferidas na versão 2 do artigo (jan. 2026); confirmar na versão de 2024 que está no espaço do Perplexity.
\begin{itemize}
    \item Teammate.next, parceria adaptativa e personalizada: o agente aprende com o retorno humano e com a própria exploração, e retém ao longo do tempo o contexto recorrente do projeto e as preferências do desenvolvedor;
    \item IDE.next, desenvolvimento orientado a intenção e a conversa, em que a intenção pode ser descrita de forma informal, por pseudocódigo, por esboços de interface ou por dados de exemplo;
    \item Compiler.next, síntese de código com múltiplos objetivos, que busca o melhor compromisso entre acurácia, latência e custo; e
    \item Runtime.next, execução ciente de SLA, com roteamento por prioridade e observabilidade para alocar recursos, estendida à computação de borda.
\end{itemize}
Nesse paradigma, o humano formula e refina as intenções, desce ao código quando precisa, e encerra o ciclo quando o software o satisfaz.

Há evidência empírica de que agentes autônomos já participam do desenvolvimento real. O estudo AIDev, de \cite{AIDEV2025}, analisou mais de 456 mil \textit{pull requests} submetidos por agentes autônomos em 61 mil repositórios do GitHub. Agentes como o Codex \citep{Codex2025} fecharam \textit{pull requests} mais rapidamente que humanos, mas suas taxas de aceitação ficaram abaixo das contribuições humanas em funcionalidades complexas e correção de defeitos \citep{AIDEV2025}. O estudo mostra o alcance e os limites dos agentes em produção, mas não controla o modo como cada agente é organizado.

Este trabalho se situa dentro da ES 3.0. Ele não compara o desenvolvedor com copiloto ao agente autônomo; os dois braços do experimento são agentes autônomos. A pergunta é interna ao paradigma: como organizar os agentes, e se a organização em papéis muda o resultado.

\subsection{Agentes Inteligentes: Fundamentos}
\label{ssec:agentes}

O ponto de partida para definir agente é o trabalho de \cite{IntelligentAgents1995}, que distingue uma noção fraca e uma noção forte de agência. Na noção fraca, um agente é um sistema computacional, em geral de software, com quatro propriedades: \textbf{autonomia}, pois opera sem intervenção humana direta e controla suas ações e seu estado interno; \textbf{habilidade social}, pois interage com outros agentes e com humanos por meio de alguma linguagem de comunicação; \textbf{reatividade}, pois percebe o ambiente e responde a mudanças nele; e \textbf{proatividade}, pois exibe comportamento orientado a objetivos por iniciativa própria. Na noção forte, o agente é descrito também por atitudes mentais, das quais o modelo \textit{Belief-Desire-Intention} (BDI), desenvolvido por \cite{Intention1987} e formalizado por \cite{BDIAgents1995}, é o mais adotado nas linguagens de programação de agentes \citep{ReviewMAS2021}.

Além dessas propriedades, os agentes atuais incorporam capacidades que este trabalho materializa em papéis. O \textbf{planejamento} decompõe objetivos complexos em sequências de ações executáveis; no braço orquestrado, ele corresponde ao planejador, que divide a tarefa em sub-tarefas. A \textbf{memória} abrange o contexto da tarefa em curso e o histórico que informa decisões futuras; no braço orquestrado, ela corresponde à memória de trabalho que o developer recebe entre ciclos. A \textbf{tomada de decisão} em ambiente parcialmente observável exige escolhas com informação incompleta; no braço orquestrado, a decisão de entregar passa pelo revisor, que aprova ou reprova o resultado.

Os LLMs mudaram a prática dos agentes. Os agentes clássicos usavam técnicas simbólicas e exigiam representação explícita do conhecimento do domínio. Os agentes baseados em LLMs usam o modelo de linguagem como núcleo e o integram a ferramentas externas, a mecanismos de planejamento e a módulos de memória. Com essa integração, um único agente navega por um repositório, lê e altera código e executa testes a partir de uma intenção expressa em linguagem natural.

\subsection{Sistemas Multiagentes: Coordenação, Cooperação e Arquiteturas}
\label{ssec:sma}

A dificuldade de um único agente diante de problemas que exigem conhecimentos de vários domínios motivou os Sistemas Multiagentes (SMA). Um SMA é um sistema computacional composto por múltiplos agentes que interagem, cada um com capacidades, objetivos e conhecimentos próprios, para resolver problemas difíceis ou impossíveis para um agente isolado. \cite{SurveyLLMs2025} atribuem aos SMA quatro propriedades: \textbf{flexibilidade e escalabilidade}, pois agentes podem ser adicionados, removidos e modificados; \textbf{robustez e confiabilidade}, pois o controle descentralizado tolera a falha de componentes; \textbf{auto-organização e coordenação}, pois os agentes dividem o trabalho e coordenam decisões por regras de comportamento; e operação em tempo real, com respostas situacionais sem supervisão humana.

A convergência entre LLMs e SMA deu origem aos \textit{LLM-based Multi-Agent Systems}, nos quais um LLM atua como núcleo de cada agente dentro de um sistema maior. O \textit{survey} de \cite{SurveyLLMs2025} sobre mecanismos de colaboração nesses sistemas conclui que eles superam abordagens de agente único quando os mecanismos de colaboração são bem projetados. O mesmo trabalho registra uma ressalva, apoiada em um estudo sobre tarefas de raciocínio: um sistema multiagente com desenho subótimo dos seus canais de colaboração competitiva pode ser superado por um agente único com prompts fortes, inclusive com demonstrações \citep{SurveyLLMs2025}. \textbf{Essa ressalva é a hipótese rival que o experimento deste trabalho testa}: com o mesmo modelo, as mesmas ferramentas, o mesmo prompt e o mesmo orçamento, a organização em papéis pode não produzir, por si, resultado diferente do agente único.

\cite{SurveyLLMs2025} caracterizam os mecanismos de colaboração em cinco eixos:
\begin{enumerate}
    \item \textbf{Atores}: os agentes envolvidos e os papéis que assumem.
    \item \textbf{Tipos de colaboração}: \textit{cooperação}, quando os agentes buscam objetivos compartilhados, \textit{competição}, quando buscam objetivos conflitantes, ou \textit{coopetição}, quando combinam os dois.
    \item \textbf{Estratégias de colaboração}: baseadas em regras, em papéis ou em modelos.
    \item \textbf{Estruturas de comunicação}: \textit{centralizada}, com um coordenador, \textit{descentralizada e distribuída}, com comunicação direta entre agentes, ou \textit{hierárquica}, com níveis de responsabilidade.
    \item \textbf{Coordenação e orquestração}: arquitetura \textit{estática}, definida previamente, ou \textit{dinâmica}, adaptada ao ambiente e à tarefa.
\end{enumerate}

Esses eixos posicionam o braço orquestrado deste trabalho. Os atores são quatro papéis sobre o mesmo modelo de linguagem. O tipo é de cooperação: os papéis buscam o mesmo objetivo, a correção da tarefa. A estratégia é baseada em papéis. A estrutura de comunicação é centralizada por um grafo, que define quem recebe o quê. A coordenação é estática: a sequência refinador de intenção, planejador, developer e revisor é fixa, com ciclos de volta do revisor ao developer e, quando o plano se esgota, ao planejador. A orquestração é feita por LangGraph.

A especialização de agentes é o ponto em que os SMA tocam a Engenharia de Software. O MetaGPT \citep{MetaGPT2023} atribui a instâncias diferentes de LLMs os papéis de gerente de produto, arquiteto, gerente de projeto, engenheiro e engenheiro de QA; o ChatDev \citep{ChatDev2023} usa CEO, CTO, programador, revisor e testador. Os dois relatam artefatos de melhor qualidade que os de abordagens de agente único. Suas comparações com agente único não isolam a topologia, porque o sistema comparado difere também em prompts, ferramentas e fluxo. As ablações internas mantêm o modelo, mas o ChatDev remove só o texto de papel dos prompts, com o mesmo número de agentes e o mesmo fluxo, e o MetaGPT acrescenta, com cada papel, etapas e artefatos do seu procedimento operacional. São os trabalhos que este estudo confronta, conforme a Subseção~\ref{ssec:correlatos}.

\subsection{Orquestração de Agentes: Grafos, Estado e \textit{Workflows}}
\label{ssec:orquestracao}

A orquestração é o conjunto de processos que governa como agentes são criados, ativados, coordenados e monitorados. Ela inclui a definição do \textit{workflow}, o gerenciamento do estado compartilhado, o roteamento das tarefas e a observabilidade da execução. Arquiteturas atuais representam o fluxo como um grafo direcionado: os nós são agentes ou funções, e as arestas são dependências e fluxos de controle. O LangGraph \citep{LangGraph2025} segue esse modelo. Ele permite execução sequencial e paralela, aceita ciclos condicionais, como os de revisão e refinamento, e persiste o estado entre as etapas. A representação em grafo também torna explícitas as dependências entre agentes e a sequência das decisões.

O gerenciamento de contexto é um dos problemas centrais da orquestração. O CodeAir \citep{CodeAir2025} trata o problema com limpeza e reinjeção de contexto entre fases: o contexto relevante volta a cada fase por meio de prompts estruturados e de artefatos persistidos. O autor relata, em uso próprio em produção e sem comparação controlada com um agente único, redução de 96\% no tempo de desenvolvimento de funcionalidades. O CodeAir é o antecedente direto da memória de trabalho do developer no Uranium: um registro determinístico, entregue entre ciclos, que substitui o histórico completo da conversa. A discussão da Seção~\ref{resultados} retoma essa relação com a medida do contexto da última chamada nos dois braços.
% Conferir no PDF do CodeAir: sessão limpa por fase e artefatos em git (não verificados nesta revisão).

A observabilidade é outra dimensão da orquestração. Num sistema com vários agentes, é preciso registrar qual agente executou qual ação, com quais entradas e saídas, em quanto tempo e a que custo. Esse registro serve à depuração e à análise científica. O Uranium usa telemetria própria: um arquivo de eventos por run (\texttt{events.jsonl}), com cada chamada ao modelo e cada uso de ferramenta; um manifesto (\texttt{manifest.json}), com toda a configuração do run; e um arquivo de resultado (\texttt{result.json}), com o julgamento do oráculo, o detector de trapaça, o custo e as medidas de qualidade. Essa telemetria permite calcular custo por tarefa e atribuir cada evento a um papel.

\subsection{Capacidades de Agentes Baseados em LLMs}
\label{ssec:capacidades}

Um agente baseado em LLM integra capacidades que permitem ao modelo agir no ambiente, planejar e manter contexto ao longo do tempo; \cite{Harmony2025} organizam essas capacidades em arquitetura, memória, planejamento e tecnologias de apoio.

\subsubsection{Utilização de Ferramentas}

A utilização de ferramentas (\textit{tool calling} ou \textit{function calling}) permite ao agente invocar recursos externos, como interpretadores de código, mecanismos de busca e sistemas de arquivos. Com essa capacidade, o LLM executa ações, verifica os resultados e usa essa informação na decisão seguinte. No Uranium, os dois braços recebem as mesmas sete ferramentas, na mesma ordem: listar, ler, buscar, escrever, substituir, \texttt{run\_python} e \texttt{run\_tests}.

\subsubsection{Planejamento}

O planejamento é a capacidade de decompor objetivos de alto nível em sequências de ações executáveis. \cite{Huang2024Planning} organizam as abordagens de planejamento de agentes baseados em LLMs em cinco linhas: decomposição de tarefas, seleção entre planos, apoio de módulos externos, reflexão e refinamento, e planejamento com memória. O \textit{Chain-of-Thought} faz o modelo gerar passos intermediários de raciocínio antes da resposta final \citep{Wei2022CoT}. O \textit{Tree of Thoughts} generaliza essa ideia: o modelo explora vários caminhos de raciocínio, avalia as próprias escolhas e volta atrás quando necessário \citep{Yao2023ToT}. O \textit{ReAct} intercala passos de raciocínio com ações no ambiente, de modo que o resultado de cada ação informa o passo seguinte \citep{Yao2023ReAct}. Em arquiteturas com mais de um componente, o planejamento ocorre em mais de um nível: um planejador decompõe o objetivo em sub-tarefas, que são delegadas a executores. O HuggingGPT segue esse esquema, com um LLM que planeja as sub-tarefas e as distribui a modelos especialistas \citep{Shen2023HuggingGPT}; o MetaGPT o aplica ao desenvolvimento de software, com papéis especializados \citep{MetaGPT2023}. O braço orquestrado do Uranium segue o mesmo esquema, com um planejador só de leitura e um developer que trata uma sub-tarefa por vez.

\subsubsection{Memória}

A memória de um agente se organiza em camadas. \cite{AINativeMemory2024} comparam o contexto do LLM à memória de curto prazo e defendem que as conclusões importantes sejam transferidas a uma memória de longo prazo, fora do contexto; \cite{Harmony2025} distinguem igualmente memória de curto e de longo prazo, com armazenamento externo para a segunda. Os primeiros propõem uma memória \textit{AI-native} que armazena e organiza conclusões derivadas do raciocínio, e não só dados brutos, e que se atualiza quando uma memória nova corrige a anterior. A janela de contexto, porém, não é neutra: o desempenho dos modelos cai à medida que o contexto cresce, mesmo quando a informação necessária está presente \citep{Levy2024SameTask, Liu2024LostInTheMiddle, Hsieh2024RULER, Modarressi2025NoLiMa}. Esse fenômeno sustenta o desenho da memória de trabalho do braço orquestrado, que mantém o contexto curto ao substituir o histórico por um registro determinístico de arquivos lidos, buscas, escritas, suítes rodadas, nota de ciclo e \textit{diff} atual.

\textit{Retrieval-Augmented Generation} (RAG) e \textit{Model Context Protocol} (MCP) ficaram fora do desenho do Uranium, nos dois braços.

\subsubsection{Síntese}

Essas capacidades costumam ser usadas para distinguir agentes de copilotos. Neste trabalho, a distinção é outra: os dois braços são agentes, com as mesmas ferramentas, o mesmo modelo e o mesmo orçamento. O agente único concentra planejamento, execução e verificação numa só conversa, e um operador determinístico roda a suíte quando ele declara ter terminado. O braço orquestrado distribui essas capacidades em papéis: refinador de intenção, planejador, developer e revisor, mais um gerador de plano de testes ao fim. A diferença entre os braços é a topologia.

\subsection{Programação Baseada em Agentes}
\label{ssec:programacao}

A programação de sistemas multiagentes precede os LLMs. A revisão de \cite{ReviewMAS2021} cataloga as linguagens de programação orientadas a agentes (APLs) e aponta, entre os obstáculos à sua adoção, a distância entre os conceitos dessas linguagens e os paradigmas de programação mais usados. Os LLMs reduzem essa distância: com um modelo de linguagem como núcleo, o comportamento do agente passa a ser especificado em linguagem natural. Uma nova geração de \textit{frameworks}, como LangChain \citep{Chase2022LangChain}, AutoGen \citep{Wu2024AutoGen} e CrewAI \citep{Moura2023CrewAI}, oferece primitivas para compor modelos, ferramentas e memória em \textit{workflows}. O LangGraph \citep{LangGraph2025} representa o \textit{workflow} como grafo direcionado, com estado compartilhado, execução condicional e ciclos. Este trabalho escolheu o LangGraph porque o grafo explicita a topologia, que é o fator do experimento, e porque os ciclos condicionais expressam a volta do revisor ao developer.

\subsection{Integridade e Rastreabilidade da Avaliação}
\label{ssec:integridade}

À medida que agentes ganham autonomia, a confiança em seus resultados depende de mecanismos verificáveis; \cite{Trustworthy2023} colocam como desafio central dos sistemas multiagentes centrados no humano projetar sistemas autônomos que sejam confiáveis e nos quais se possa confiar, e propõem avaliar a confiabilidade de um agente pelo seu comportamento observável. Na avaliação de agentes de código, essa exigência se traduz em julgar o resultado com evidência que o agente não controla. O SWE-bench \citep{Jimenez2024SWEbench} consolidou esse desenho: cada tarefa vem de um defeito real, e o julgamento usa testes escritos pelo projeto, que devem passar de falha a sucesso (\textit{fail-to-pass}) sem quebrar os que já passavam (\textit{pass-to-pass}); o SWE-bench Verified restringiu o conjunto a tarefas validadas por humanos \citep{OpenAI2024SWEbenchVerified}. Trabalhos posteriores apontam como ameaças a esse tipo de avaliação a contaminação dos dados de treino \citep{Jain2025LiveCodeBench}, testes insuficientes que aprovam patches errados \citep{Liu2023EvalPlus, Yu2025UTBoost, Wang2026SolvedIssues} e a manipulação dos próprios testes pelo agente (\textit{gaming}) \citep{Zhong2026ImpossibleBench}.

O Uranium implementa os mecanismos que essa literatura pede. O oráculo julga cada run, depois de encerrado, com testes que o agente nunca viu, em três verificações: \textit{fail-to-pass}, \textit{pass-to-pass} e regressão. A área de testes é protegida, e toda execução de código do agente ocorre numa caixa de areia sem rede. O workspace nasce só com o commit-base, sem histórico \textit{git}, o que impede o acesso a código do futuro do projeto. O detector de trapaça aplica oito regras estáticas sobre o patch. O manifesto e o arquivo de eventos registram a configuração e a execução de cada run, e a semente é a mesma nos dois braços de cada par. A Seção~\ref{desenvolvimento} detalha esses mecanismos.

A avaliação de experimentos pareados com poucas unidades pede testes exatos e análise de poder. Este trabalho usa o teste do sinal exato bilateral por tarefa, o McNemar exato por par e o Wilcoxon exato, além de poder estimado por simulação, seguindo guias de experimentação em Engenharia de Software \citep{Wohlin2012Experimentation, Arcuri2014Hitchhiker}. Os testes são os de \cite{Dixon1946SignTest}, \cite{McNemar1947} e \cite{Wilcoxon1945}, nas versões exatas \citep{Hollander2014Nonparametric, Fagerland2013McNemar}.

\cite{GenAIAgenda2023} observam que o impacto das ferramentas de IA generativa tem sido discutido sobretudo pela produtividade, e que elas reabrem as preocupações clássicas sobre produtividade e qualidade.
% Supõe que GenAIAgenda2023 é Nguyen-Duc et al.; a frase anterior ("complementar desempenho com qualidade e manutenibilidade") não está nesse artigo.
Este trabalho mede sete métricas, com hierarquia fixada antes da coleta: resolução; tokens; tempo de parede; consistência; variação da complexidade ciclomática; conclusão sem corte; e violações do detector. O custo em dólares é medido à parte. As dimensões do projeto original se mapeiam assim: produtividade, sem humanos no laço, é a eficiência do sistema e vira o tempo de parede; autonomia vira conclusão sem corte; qualidade vira a variação da complexidade ciclomática; custo vira tokens, com dólares e chamadas reportados; e rastreabilidade deixa de ser métrica e vira portão de admissão do dado, com provedor pinado em toda chamada, rede isolada na avaliação e trilha de eventos válida. Das métricas de qualidade do código, só a complexidade ciclomática entrou. O Índice de Manutenibilidade é calculado, mas saturou em zero nos módulos alterados e foi excluído. O \textit{Mutation Score} foi excluído porque mede a força de uma suíte de testes, e nenhum braço escreve testes, já que a área de testes é protegida; ele mediria a suíte do repositório, idêntica nos dois braços.

\subsection{Trabalhos Correlatos}
\label{ssec:correlatos}

Os trabalhos correlatos se agrupam em quatro linhas: sistemas multiagentes com papéis para desenvolvimento de software; estudos empíricos de agentes em repositórios reais; avaliação de agentes de código por testes escondidos; e comparações entre topologias, falhas e custo de sistemas multiagentes.

\textbf{Sistemas multiagentes com papéis.} O MetaGPT \citep{MetaGPT2023} e o ChatDev \citep{ChatDev2023} organizam LLMs em papéis inspirados em equipes humanas e relatam ganhos sobre abordagens de agente único. O CodeAir \citep{CodeAir2025} organiza o desenvolvimento em fases com limpeza e reinjeção de contexto e relata, sem comparação controlada, redução de tempo em uso próprio em produção. Os três avaliam o sistema inteiro. O modelo, os prompts e as ferramentas do sistema proposto diferem dos usados na comparação, e o efeito da organização em papéis fica confundido com essas diferenças. Este trabalho mantém essas variáveis constantes e faz variar só a topologia.

\textbf{Mecanismos de colaboração.} O \textit{survey} de \cite{SurveyLLMs2025} oferece a caracterização em cinco eixos e a ressalva de que um agente único com prompts fortes pode superar sistemas multiagentes com canais de colaboração mal desenhados. Este trabalho testa essa ressalva num cenário controlado, com uma configuração fixa do braço orquestrado.

\textbf{Agentes em repositórios reais.} O AIDev \citep{AIDEV2025} mostra que agentes autônomos já submetem \textit{pull requests} em escala, com aceitação menor que a humana em tarefas complexas. O estudo é observacional e não controla a organização interna dos agentes. Este trabalho usa defeitos e funcionalidades reais minerados de projetos Python abertos, mas em ambiente controlado.

\textbf{Avaliação por testes escondidos.} O SWE-bench e seus derivados verificados definem o julgamento por \textit{fail-to-pass} e \textit{pass-to-pass} que o oráculo do Uranium adota \citep{Jimenez2024SWEbench, OpenAI2024SWEbenchVerified}. As críticas de contaminação, de testes insuficientes e de \textit{gaming} motivam o detector de trapaça e o workspace sem histórico \citep{Jain2025LiveCodeBench, Yu2025UTBoost, Zhong2026ImpossibleBench}.

\textbf{Agente único contra multiagente com controle de variáveis.} Estudos que comparam topologias com o modelo constante são os correlatos mais diretos deste trabalho. \cite{Dong2024SelfCollaboration} fazem analista, programador e testador sobre o mesmo ChatGPT superar o ChatGPT direto em geração de código; \cite{Xia2025Agentless} mostram que uma esteira fixa de localização, reparo e validação, sem agente autônomo, compete com agentes no SWE-bench Lite a custo menor; e \cite{Kim2025ScalingAgentSystems}, em preprint, comparam um agente único a quatro arquiteturas multiagente com prompts, ferramentas e orçamento idênticos e encontram ganhos limitados no SWE-bench Verified quando o agente único já é forte. Nenhum deles mantém, ao mesmo tempo, modelo, ferramentas, prompt, orçamento e semente constantes entre um agente único e uma orquestração de papéis, em defeitos reais julgados por testes escondidos, com custo, tempo, consistência e integridade medidos por par. Essa combinação é a lacuna que este estudo trata.

\textbf{Falhas de sistemas multiagentes.} \cite{Cemri2025MAST} classificam as falhas de sistemas multiagentes com LLMs em problemas de desenho do sistema, de desalinhamento entre agentes e de verificação da tarefa. Essas categorias dão nome a mecanismos que a Seção~\ref{resultados} analisa, como o revisor que reprova um patch certo e o plano que serializa defeitos independentes.

\textbf{Custo e eficiência de agentes.} Trabalhos que medem tokens e dólares por tarefa resolvida permitem situar o custo do braço orquestrado \citep{Kapoor2025AgentsThatMatter, Kapoor2026HAL}. A Seção~\ref{resultados} apresenta esse custo, com o recorte da carga inteira marcado como análise exploratória.

Em síntese, a literatura propõe sistemas multiagentes com papéis e os avalia como sistemas completos, enquanto a avaliação de agentes de código já dispõe de oráculos por testes escondidos. Falta o experimento que combine os dois: um oráculo de testes escondidos e uma comparação em que só a topologia varie. Este trabalho ocupa esse espaço.
